Agents that learn from other agents must estimate how far to trust each source.
We study a population of $N$ neural-network learners under a Bayesian online rule
that jointly updates each agent's weights and an explicit, dynamical reliability
estimate for every other agent. Into this population we introduce a
\emph{representation error}: a fraction $f_d$ of agents extend their encoding of
each message with a feature of strength $d$ that carries no information about the
task and depends only on the sending agent's group label. Individually this is a
mild misspecification, and against a randomly assigned label it does nothing but
degrade learning. Against the group label it does not stay mild: above a sharp
boundary in $(d, f_d)$ the population splits into mutually distrustful groups
aligned with the label. We map the full phase diagram and identify two distinct
discriminatory regimes, one in which the split still tracks learned
representations and one in which trust tracks the label almost independently of
them. Ablations isolate the cause. Freezing the reliability sector, or removing
the sign reversal by which distrusted sources are anti-learned, eliminates the
transition; Hebbian and Perceptron updates under the same perturbation do not
produce it at all. Finally, task diversity $\alpha = P/K$ controls the
\emph{order} in which structure appears: with few tasks, trust structure forms
first and drags representations behind it, and with many the reverse. The
amplification is therefore a property of reliability-weighted learning rather than
of the architecture, and it is strongest in exactly the regime such weighting is
usually adopted for.
